\worksheettitle
  {Домашняя работа}
  {\modulemark{M07-H} Условия между цифрами}

\studentnameblock

\begin{notebox}[Маршрут]
  В каждом задании сначала обозначьте цифру десятков через \(a\), цифру
  единиц через \(b\). Запишите условие математически и выберите цифру, по
  которой удобнее вычислять вторую. Проверьте по порядку все значения
  выбранной цифры и оставьте только пары, где
  \(1\leq a\leq9\), \(0\leq b\leq9\).
\end{notebox}

\begin{practicebox}[1. Запишите условие математически]
  \begin{enumerate}[label=\textbf{\arabic*.}]
    \item Цифра единиц на \(3\) больше цифры десятков:
      \answerblank[38mm].
    \item Цифра десятков на \(2\) больше цифры единиц:
      \answerblank[38mm].
    \item Цифра единиц равна удвоенной цифре десятков:
      \answerblank[38mm].
    \item Сумма цифр равна \(11\): \answerblank[38mm].
  \end{enumerate}
\end{practicebox}

\begin{practicebox}[2. Найдите полный список]
  Цифра единиц на \(5\) больше цифры десятков.

  Запись: \answerblank[32mm]. Перебираю: \answerblank[18mm].

  Список: \answerblank[100mm].

  Какие значения проверены? \answerblank[65mm].
\end{practicebox}

\begin{practicebox}[3. Найдите полный список]
  Цифра единиц на \(4\) меньше цифры десятков.

  Запись: \answerblank[32mm]. Перебираю: \answerblank[18mm].

  Список: \answerblank[100mm].

  Почему список полный? \writinglines{1}
\end{practicebox}

\newpage
\begin{practicebox}[4. Кратная зависимость]
  Найдите все двузначные числа, в которых цифра десятков в \(3\) раза больше
  цифры единиц.

  Запись: \answerblank[32mm]. Перебираю: \answerblank[18mm].

  Список: \answerblank[100mm].

  Почему других чисел нет? \writinglines{1}
\end{practicebox}

\begin{practicebox}[5. Сумма цифр]
  Найдите все двузначные числа с суммой цифр \(12\).

  Запись: \answerblank[32mm]. Перебираю: \answerblank[18mm].

  Список: \answerblank[100mm].

  Количество чисел: \answerblank[15mm].
\end{practicebox}

\begin{warningbox}[6. Найдите и исправьте ошибки]
  \begin{enumerate}[label=\textbf{\arabic*.}]
    \item Для \(b=a+5\) после числа \(49\) включили пару \((a,b)=(5,10)\).
    \item Для суммы цифр \(3\) записали \(12,21\), но не записали \(30\).
    \item Для \(a=3b\) записали \(13,26,39\).
  \end{enumerate}
  Для каждой записи назовите причину ошибки и дайте правильный ответ.
  \writinglines{3}
\end{warningbox}

\newpage
\begin{practicebox}[7. Номер участника]
  Участнику соревнования присвоили двузначный номер меньше \(60\). Сумма цифр
  номера равна \(9\). Какие номера могли присвоить участнику?
  \writinglines{2}
\end{practicebox}

\begin{practicebox}[8. Обратная задача]
  Числа \(15,26,37,48,59\) получены по одному условию между цифрами.
  \begin{enumerate}[label=\textbf{\arabic*.}]
    \item Сформулируйте условие словами. \writinglines{1}
    \item Запишите его равенством: \answerblank[38mm].
    \item Объясните, почему пара \((a,b)=(6,10)\) не даёт следующего числа.
      \writinglines{1}
  \end{enumerate}
\end{practicebox}

\begin{checkpointbox}[9. Проверка перед сдачей]
  Выберите одно задание 2--5 и докажите полноту списка: назовите, какая цифра
  изменялась по порядку, перечислите все проверенные значения, объясните, как
  вычислялась вторая цифра и почему все отброшенные кандидаты недопустимы.
  \writinglines{3}
\end{checkpointbox}
